\documentclass[12pt,letterpaper]{letter}
\usepackage[T1]{fontenc}
\usepackage{newtxtext,newtxmath}
\usepackage[margin=1in]{geometry}
\usepackage{microtype}
\usepackage[hidelinks]{hyperref}
\setlength{\parskip}{0.7em}
\setlength{\parindent}{0pt}
\signature{Lily Luo\\{}[AUTHOR TO COMPLETE: affiliation]\\{}[AUTHOR TO COMPLETE: email]}
\address{[AUTHOR TO COMPLETE: postal address]}
\begin{document}
\begin{letter}{Editors\\The Review of Economics and Statistics}
\opening{Dear Editors,}

I submit ``What Is AI Exposure? Measurement Architecture and Statement-Specific Robustness in Early-Career Employment'' for consideration at \emph{The Review of Economics and Statistics}. The paper asks how economic conclusions should be evaluated when several defensible constructions exist for the same named treatment. It follows six occupational AI-exposure architectures from their technological and occupational primitives through taxonomy, empirical support, employment-stock contrasts, and worker-mobility rankings.

Three findings organize the paper. First, the architectures differ materially in occupational content, rankings, and effective identifying support. Second, all six produce negative high-versus-low early-career employment-stock point estimates on literal common support. Third, that directional agreement does not make the measures interchangeable: mobility labels frequently conflict, and an exploratory family decomposition shows conditional asymmetry that is sensitive to the distinctive direct-exposure measure.

The contribution extends beyond this application. Tail comparisons, continuous regressions, and mobility claims consume different representations of a constructed treatment. The paper develops and demonstrates the principle that robustness is statement-specific, connecting upstream measurement choices to the exact economic statements they can support.

The empirical boundary is explicit. Exposure measures potential occupational susceptibility rather than observed adoption, and the employment estimates are associations rather than causal effects. The design and confirmatory analyses were frozen before protected outcomes were opened; subsequent exploratory analyses retain their status labels and reproducibility receipts. [AUTHOR TO COMPLETE: exclusivity, conflicts, and prior-publication declaration.]

\closing{Sincerely,}
\end{letter}
\end{document}
